\documentclass[../main.tex]{subfiles}
\begin{document}
%==========================================TIÊU ĐỀ TẬP========================================
\begin{table}[h]
  \begin{adjustwidth}{-1cm}{}
    \begin{tabular}{>{\centering\arraybackslash}p{6cm}|>{\raggedright\arraybackslash}p{12.5cm}}
      \multirow{5}{6cm}
      % Cột bên trái: ÉPISODE và số tập
      {\centering
        {\\[-40pt]
          \flaregothic\fontsize{20pt}{20pt}\selectfont
          \color{eptype}ÉPISODE\color{black}\\
          \burbank\fontsize{72pt}{72pt}\selectfont
          \color{epnum}119
          \\[10pt]
          \hrule
          \vspace{8pt}
          \flaregothic\fontsize{20pt}{20pt}\selectfont
          \color{eptype}SPÉCIAL\color{black}\\
          \burbank\fontsize{72pt}{72pt}\selectfont
          \color{epnum}XXIV\color{black}\\[6pt]
          \cafeteria\fontsize{20pt}{20pt}\selectfont\color{epnum}(S-24)\color{black}
          \\[18pt]
          % \flaregothic\fontsize{14pt}{14pt}\selectfont
          % \color{epattr}BIENVENUE, \uline{\color{Banpire} Mel} !
          \flaregothic\fontsize{18pt}{18pt}\selectfont
          \color{gray} COMPILATION\\
          \color{black}
        }
      }
       &
      % Dòng thứ nhất: tên tập tiếng Nhật
      {\fotskip\fontsize{18pt}{18pt}\selectfont \ruby{終}{しゅう}\ruby{焉}{えん}のトリガーになった\ruby{瞬}{しゅん}\ruby{間}{かん}\topten
      } \\[6pt]
      % Dòng thứ hai: tên tập tiếng Anh
       & {\cafeteria\fontsize{18pt}{18pt}\selectfont The Beginnings of the End}                                       \\[4pt]
      % Dòng thứ ba: tên tập tiếng Việt
       & {\vnmsans\fontsize{18pt}{18pt}\selectfont Top 10 - Các tập "tận thế" nhất}                                   \\[8pt]
      % Dòng thứ tư: Nhân vật xuất hiện
       & {\caviar{\fontsize{14pt}{14pt}\selectfont Nhân vật: [\emp{kuon1} \emp{achan}] ({\abv Truyện phụ}) \emp{sora}}} \\[8pt]
      % Dòng cuối cùng: Thời gian diễn ra của tập (thường sẽ là ngày phát hành)
       & {\continuum\fontsize{16pt}{16pt}\selectfont Ngày phát hành: 22/08/2021}                                      \\[36pt]
    \end{tabular}
  \end{adjustwidth}
\end{table}
%=========================================TRUYỆN CHÍNH========================================
\color{black}

\vspace{60pt}

\begin{center}
  (Đây là danh sách top 10 do fan bình chọn cho những tập hỗn loạn nhất của \textit{holo no Graffiti}.)
\end{center}

Ánh nắng sớm nhẹ nhàng lan tỏa khắp thung lũng xanh ngát, tiếng chim hót líu lo vang vọng giữa khung cảnh núi đồi yên ả. Sora đứng giữa trung tâm hình ảnh, mặc bộ trang phục biểu diễn rực rỡ với nụ cười nở rộng trên gương mặt. Gió thoảng nhẹ thổi qua, làm làn tóc cô khẽ bay bay và mang theo hương thơm từ rừng thông phía sau lưng.

Vừa đưa tay lên vừa nháy mắt tinh nghịch, cô vui vẻ chào đón khán giả: ``Tóm tắt câu chuyện vừa rồi!''

Nhưng lời vừa cất lên xong, bầu trời trong xanh bỗng chốc chuyển sang màu đỏ thẫm như bị lửa thiêu đốt. Tiếng rít của ngọn lửa và kim loại va chạm kêu lên inh tai. Cảnh vật phía sau Sora lập tức biến thành địa ngục trần gian: những tòa nhà cháy rụi, các mảng tường liên tục nổ tung, mặt đất nứt vỡ toang và bụi khói cuồn cuộn bốc lên như những cột lửa. Những tia sét đỏ từ sâu dưới lòng đất phóng thẳng lên trời, tạo nên thứ ánh sáng dữ dội đến chói lóa mắt.

\begin{figure}[H]
  \centering
  \begin{subfigure}[t]{0.45\textwidth}
    \centering
    \includegraphics[width=8cm]{../../../../Image/Book9/s-22a.png}
  \end{subfigure}
  \quad
  \begin{subfigure}[t]{0.45\textwidth}
    \centering
    \includegraphics[width=8cm]{../../../../Image/Book9/s-22b.png}
  \end{subfigure}
  \caption{Hai thái cực đối lập nhau trong phần mở đầu.}
\end{figure}

Sora vẫn đứng nguyên vị trí, nhưng lần này gương mặt cô mang vẻ nghiêm trang đến lạ thường. Nụ cười đã tan biến, thay vào đó là ánh mắt đầy hoang mang và giọng nói trầm lắng, chậm rãi: ``Bọn mình đã phạm phải sai lầm gì đâu? Tại sao chuyện này lại có thể xảy ra?''

Cô quay đầu nhìn về phía sau, nơi một tòa nhà đang nghiêng ngả trước khi sụp đổ như bị hút vào khoảng không vô tận. Một chùm đèn sân khấu bật sáng, dõi theo từng bước chân của Sora khi cô tiếp tục kể chuyện: ``Mình không thể cứ thế mà chấp nhận số phận bi thảm này được.''

Rồi cô quay mặt về phía khán giả, hay chính xác hơn là người đang xem đoạn băng ghi hình, nhìn thẳng vào ống kính với ánh mắt quyết đoán: ``Mình phải quay trở lại những sự kiện trong quá khứ để tìm ra gốc rễ của mọi chuyện.''

Một khoảng lặng ngắn ngủi trôi qua. Bất ngờ, gương mặt cô lại bừng sáng như có công tắc được bật lên, đôi mắt cong lên theo nụ cười và giọng nói vui vẻ trở lại như mọi khi: ``Chúng ta cùng xem lại nào!''

Cô giơ một tay lên chỉ thẳng vào khán giả, giọng nói vang vọng như người dẫn chương trình truyền hình:

\begin{center}
  \textbf{``Top 10 cảnh của \textit{hologra} có khả năng tạo thành ngày tận thế' xin được bắt đầu ngay bây giờ!''}
\end{center}

Và tất cả những lời đó đều được nói ra giữa khung cảnh đang cháy rụi, với đống đổ nát còn âm ỉ lửa sau lưng, vài mảnh vỡ pháo hoa rơi rào rào từ bầu trời và những bóng người không rõ mặt đang chạy tán loạn. Sora vẫn giữ nguyên nụ cười tươi rói, một niềm vui kỳ lạ giữa cơn đại họa, như thể tất cả những điều này chỉ là một tiết mục nhẹ nhàng trong buổi sáng bình thường.

\il

\begin{center}
  \uline{(Phần bình luận của Sora sẽ được viết dưới dạng thoại.)}
\end{center}

\pov{Tokino Sora pha lẫn vào ngôi kể thứ ba}

\bxh{10}{Con tàu của Suisei!}{109}{Chín}    % 9-1

[\dots] Đúng lúc đó, Kanata vừa bước vào văn phòng. Nhìn thấy cảnh tượng lạ lùng trước mắt, cô đứng sững người một lát, sau đó quay sang hỏi người quản lý đang ngồi xoay ghế đối diện: ``Tàu\dots\ Suisei? Cái này là chuyện gì vậy, Kuon-senpai?''

Kuon nhún vai đáp lời cô nàng Thiên sứ: ``Anh cũng không rõ lắm. Bỗng dưng em ấy muốn làm tàu, thế là tự cắt dán bìa các-tông thành cả một con tàu hỏa thế này. Đã chạy vòng quanh phòng như vậy được khoảng 15 phút rồi đấy.''

Ryuuhou vẫy tay chào Kanata, giọng đầy hài hước: ``Kanata-san, mau lên tàu kẻo lỡ chuyến! Em cũng vừa định xin đi một vòng đây.''

\begin{dialogue}
  \speak{\textit{Sora}} À, cảnh mà Suisei-chan giả vờ làm đoàn tàu đây mà. Trông cô ấy cps vẻ vui đấy, nhưng đó chỉ là bề nổi thôi\dots
\end{dialogue}

Nàng Thiên sứ nhìn tiền bối với ánh mắt không thể tin nổi. Nhưng trước khi kịp nói gì, Suisei đã chạy ngay đến bên cô, đôi mắt sáng rực gọi: ``Kanata! Em đến rồi này!''

Cô nắm lấy tay Kanata và hào hứng mời: ``Em có muốn lên chuyến tàu Suisei Cáu kỉnh này không?''

Kanata giật mình ngạc nhiên: ``Đấy là chị khi bị cáu kỉnh ạ?''

Kuon bình thản tiếp lời từ phía sau: ``Có thể em không tin, nhưng đúng là thế đấy.''

Ryuuhou cũng xen vào, giọng đầy tò mò: ``Suisei-san mà cáu kỉnh thì trông thế nào nhỉ? Câu lạc bộ em có một thành viên cũng hay nổi nóng, nhưng không biết có `độ' nào bằng chị không.''

Kanata vẫn chưa hết bàng hoàng, liền hỏi thêm về con tàu hỏa kỳ lạ: ``Không phải chị thường xuyên đi diễn hát sao? Sao lại có thời gian làm cái này?''

\begin{center}
  \uline{(Có một cái đồng hồ đếm ngược ở góc dưới màn hình. Đồng hồ chỉ về không khi cảnh ở dưới xảy ra.)}
\end{center}

Nhưng ngay lập tức, cô nhận lại một ánh mắt sắc lạnh từ nàng Ca sĩ. Suisei trừng mắt, giọng nói đơn điệu mà lại đầy vẻ triết lý: ``Không có niềm vui nào trên đời này là kéo dài mãi được, đúng không?''

\img{9-1b}

``Mặt chị nhìn sợ quá!!'' Kanata run rẩy, mặt tái mét, tay bất giác bám chặt lấy cánh tay Kuon như tìm điểm tựa vững chắc.

Ryuuhou nhìn Suisei rồi cười lớn: ``Suisei-san, trông chị giống ác quỷ trong truyện vậy! Nhưng mà em thích cái vibe đó đấy.''

Kuon cũng bật cười thành tiếng: ``Đây mới chính là cái `cáu kỉnh' em ấy nói đấy. Sui-chan, đừng có dọa em ấy như thế.''

\begin{dialogue}
  \speak{\textit{Sora}} Đúng như các bạn thấy. Cơn tức giận của cậu ấy quả thực có khả năng hủy diệt cả thế giới.
\end{dialogue}

Vẻ mặt nghiêm nghị của nàng ``Sao chổi'' lập tức tan biến, trở lại vui tươi như thường. Cô vỗ nhẹ vào thân tàu bìa các-tông và mời hậu bối: ``Nói đùa thôi, em cũng lên thử đi. Cùng lên tàu với Sui-chan nào!''

Kanata do dự nhìn sang cậu Tổng, ánh mắt như muốn hỏi ``Anh có chắc là mọi thứ sẽ bình thường không?''. Ryuuhou đứng dậy, vỗ vai cô động viên: ``Kanata-san, sợ gì chứ! Em cũng định lên cùng chị nữa đây. Mà Onii, anh không lên tàu cùng mọi người à? Cơ hội hiếm có lắm đấy.''

Kuon đáp lại, cười khì và xua tay: ``Cứ lên tàu đi chơi đi. Anh đang bận việc, chắc không thể tham gia được.''

Nàng Thiên sứ miễn cưỡng trèo lên con tàu hỏa, trong khi Ryuuhou thì nhảy lên đuôi tàu, tay giơ cao như một hành khách nhiệt tình: ``Em xin một vé đi vòng quanh thế giới!''

Suisei hân hoan tuyên bố trước khi xuất phát: ``Con tàu Suisei đã sẵn sàng, khởi hành thôi!''

\begin{dialogue}
  \speak{\textit{Sora}} Mọi người đã hiểu rồi chứ? Sự kiện ở tập đó đã diễn ra như vậy đấy!
\end{dialogue}

``\textit{Vậy là chúng ta có một con tàu Suisei, với `ba mẹ con' ngồi trên đó,''}'' Kuon cười thầm trong lòng, nhìn Kanata ngồi bấp bênh và Ryuuhou đang tạo dáng như một du lịch bụi thực thụ.

``Cái quái gì mình lại dính vào thế này\dots'' Kanata nhăn mặt, không thể tin được mình lại tham gia vào trò chơi kỳ lạ này. [\dots]

\bxh{9}{Cừu uống cà phê!}{89}{Bảy}    % 7-6

[\dots] Watame vùng lên, đập tay xuống bàn với một tiếng *BỐP!*, thốt lên một câu đầy khẳng định: ``Em sẽ trở thành một cô Cừu có thể uống cà phê!''

Kuon phì cười, chống cằm nhìn cô đầy thích thú: ``Vừa nãy uống còn phun nó vào mặt Subaru-chan thế kia mà còn đòi.'' Cậu quay sang Subaru, nhếch mép cười.

Cậu Tổng phì cười, chống cằm nhìn cô đầy thích thú, giọng cậu như một lời nhắc nhở đầy hài hước: ``Vừa nãy uống còn phun nó vào mặt Subaru-chan thế kia mà còn đòi.'' Cậu quay sang nàng Vịt, nhếch mép cười với vẻ mặt đầy chọc ghẹo. ``Mà nói đến Subaru-chan thì\dots''

Ngay giây tiếp theo, một luồng sát khí từ nàng Vịt tỏa ra phía sau lưng nàng Cừu, khiến không khí bỗng chốc trở nên lạnh lẽo. Cô nàng Tomboy đặt tay lên vai nàng Cừu, giọng nói trầm xuống một cách nguy hiểm, đầy đe dọa. Cô siết nhẹ vai nàng Cừu, cất lên một giọng đáng sợ: ``\dots Xin lỗi chị đi đã.''

\begin{dialogue}
  \speak{\textit{Sora}} Đây là một cảnh mà Watame-chan sau khi phun cà phê đen vào mặt của Subaru-chan\dots
\end{dialogue}

Nàng Cừu cứng đờ người, toàn thân cô như bị đóng băng trong khoảnh khắc.

Trong khi đó, Choco giả vờ không thấy cảnh tượng đó, tiếp tục hỏi với giọng điệu đầy trêu chọc, đôi mắt cô ánh lên vẻ thích thú: ``Rồi, rồi. Em có làm được không đấy, Watame-sama?''

Nàng Cừu sà tới với vẻ mặt cực kỳ nghiêm túc, hai tay cô chống lên bàn: ``Em làm được! Em biết mà!''

\begin{center}
  \uline{(Có một cái đồng hồ đếm ngược ở góc dưới màn hình. Đồng hồ chỉ về không khi cảnh ở dưới xảy ra.)}
\end{center}

Vừa dứt lời, cô lập tức bị nàng Vịt bẻ người nhẹ ra phía sau. Cô nàng giãy nảy, cố cầm hai cánh tay đang ghì chặt cô, giọng cô đầy thống khổ: ``Uwawawwa\dots\ Tha cho emmmmmm!!''

Nàng Y tá nhấp một ngụm cà phê rồi nói tỉnh bơ, giọng cô như một nhà triết gia đang đưa ra lời khuyên: ``Vậy em nên bắt đầu bằng việc hiểu rõ xem mình đang ở đâu đi.''

\img{7-7f}

Watame vẫn đang vẫy vùng để thoát khỏi Subaru, người vẫn siết chặt như một chiếc khóa gọng kìm, không có dấu hiệu nhả ra.

\begin{dialogue}
  \speak{\textit{Sora}} Mình nghe nói là có một đấu vật siêu hạng nào đó bắt đầu sử dụng đòn này trong trận đấu của người đó rồi đấy! Nhưng mà các bạn à, đừng làm thử tại nhà nha.
\end{dialogue}

Cuối cùng, sau một hồi vật lộn, nàng Cừu thoát ra được, đứng thẳng người, đập tay xuống bàn với vẻ mặt đầy quyết tâm: ``Chúng ta cùng bắt đầu cuộc họp chiến thuật nào!''

Không khí trong phòng đột nhiên biến thành một cuộc họp quân sự.

Nàng Cừu quay sang nhìn nàng Vịt, người đang lấy khăn lau cà phê vẫn còn dính trên mặt. Cô chỉ tay thẳng về phía nàng Tomboy: ``Đi thôi nào, Subaru-senpai!''

Đàn chị của cô giật mình, tròn mắt nhìn cô nàng Cừu đầy cảnh giác, giọng cô như một tiếng kêu cứu: ``Ể\dots\ Chị không muốn đi--!?'' Nhưng trước khi cô kịp phản ứng, thì hậu bối đã kéo cô đi ra khỏi văn phòng với sức mạnh không ngờ [\dots].

\bxh{8}{Chiếc đồng hồ nổi giận!}{100}{Tám}  % S-19

[\dots]  ``PolPol tập trung nha\~{}, được chứ?'' Botan hỏi với cái miệng mở to hết cỡ, rồi đứng dậy bấm nút phóng với quyết tâm không gì lay chuyển được.

``Tao biết rõ chuyện này mà, cảm ơn nha!'' nàng Cáo Sa mạc hét lên khi bị bắn đi như một viên đạn sống, giọng reo vui như kẻ đang tận hưởng cuộc phiêu lưu mới.

\begin{dialogue}
  \speak{\textit{Sora}} Đây là cảnh mà bọn họ (chỉ bộ ba Aqua-Marine-Kuon) bị đuổi bắt bởi nhóm băng đảng trong rừng.
\end{dialogue}

``\dots Mẹ kiếp.'' Đó là những lời nói cuối cùng của Kuon trước khi ``thủ lĩnh'' của băng đảng  đáp xuống chiếc xe Mio, tạo ra một vụ nổ nhẹ khiến mọi người bị hất văng xuống đất như bắp rang bị thổi tung trong lò vi sóng.

``\dots Mẹ kiếp.'' Đó là những lời cuối cùng của Kuon trước khi kẻ đứng đầu băng đảng đáp xuống chiếc xe Mio, tạo ra một vụ nổ nhỏ khiến mọi người bị hất tung ra khỏi xe như bắp rang trong lò vi sóng; một cảnh tượng hỗn loạn tột cùng.

Chiếc Miofast hét lên thảm thiết, mất lái và đâm sầm vào chiếc đồng hồ một cách đầy định mệnh, rồi phát nổ như một bộ phim hành động kinh phí thấp, mảnh vỡ bay tứ tung như sao rơi.

\begin{center}
  \uline{(Có một cái đồng hồ đếm ngược ở góc dưới màn hình.)}
\end{center}

Tất cả mọi người, trừ nàng Thuyền trưởng và cậu Tổng, đều ngất đi hoàn toàn, thân thể nằm rải rác trên mặt đất như những con rối bị cắt dây. Chỉ còn hai người tỉnh táo, nhưng quần áo đầy bụi bẩn, mặt mày lem luốc với vẻ mệt mỏi và bất lực.

``Cái quái gì vừa xảy ra thế này...? Ái da\dots'' Kuon vừa xoa đầu vừa lồm cồm ngồi dậy, giọng như người vừa tỉnh sau cơn ác mộng. Marine đưa tay về phía đám khói, như nhân vật chính trong phim mất đi người bạn thân trong chiến tranh, giọng đầy đau thương: ``Đồng hồ ơi\dots\ Không\dots''

Kuon đứng dậy phủi bụi, gắt gỏng đầy bực bội: ``Nó bị giết chết rồi. Tuyệt cmn vời gì đâu.''

Bất chấp lời càu nhàu của cậu, nàng Thuyền trưởng vẫn rơm rớm nước mắt, nghẹn ngào nhìn chiếc đồng hồ cầm tay mà chính cô đã dọa nạt trước đó: ``Ta sẽ nhớ mãi về ngươi, đồng hồ ơi\dots''

``Anh sắp phát điên với cái nơi này rồi\dots'' cậu thở dài, thi thoảng ho vài tiếng, giọng của người đã mất hết niềm tin vào thế giới thực tại.

Đám khói trước mặt vẫn bốc lên nghi ngút, như đang che giấu một bí mật sắp được hé lộ\dots

\begin{center}
  \uline{(Đồng hồ chỉ về không khi cảnh ở dưới xảy ra.)}
\end{center}

``Hửm?'' / ``Cái quái gì thế\dots''

Và đúng lúc đó, linh hồn đầy cơn thịnh nộ của cái đồng hồ nổi dậy, theo lời kể của nàng Vu nữ: một câu chuyện về sự phục sinh và trả thù từ cõi khác.

Từ trong màn khói, một bóng dáng kỳ dị hiện ra dần. Đó là một người có thân hình giống A-chan, nhưng thay vì gương mặt bình thường, mặt cô lại là một cái đồng hồ, những chiếc kim quay với tốc độ đáng sợ. Nó lao tới như muốn đấm bay hai người, miệng lặp đi lặp lại câu: ``Trời sáng rồi mấy đứa ơi! Dậy đi thôi nào!''

\img{s-18o}

\begin{dialogue}
  \speak{\textit{Sora}} Giá như linh hồn giận dữ ấy không chiếm lấy thân xác của chiếc đồng hồ đó\dots\ Nhưng rất may cho họ đó chỉ là một giấc mơ mà thôi.
\end{dialogue}

Marine và Kuon tròn mắt, chưa kịp phản ứng thì cái mặt đồng hồ đã tung cước đấm thẳng vào đầu cả hai với lực mạnh đến mức họ bị hất bay ngược ra sau như cảnh tua chậm trong phim hành động, thân thể vẽ một đường cong trên không trước khi rơi xuống đất với tiếng động nặng nề.

Trong giây phút mất đi ý thức, họ chỉ còn nghe thấy giọng nói vang vọng của bà thần Okayu, như lời tạm biệt từ một thế giới xa xôi: ``Tạm biệt nhá, nyan\~{}''

Và rồi, màn đêm buông xuống, nhấn chìm tất cả trong sự im lặng đầy bí ẩn, cùng những câu hỏi chẳng có lời đáp. [\dots].

\bxh{7}{Tòa nhà Kuon!}{28}{Hai}   % 2-11

[\dots] Đúng lúc đó, Người dưới đất lại bất ngờ trồi lên một lần nữa, tóm chặt lấy chân cạu Tổng. Hắn nói với giọng khản đặc: ``Cậu cũng phải làm phân bón luôn cho đủ bộ!''

``\textbf{Cái gì!? Khoan đã!!!}'' Kuon hét lên kinh hãi nhưng đã bị hắn lôi tuột xuống lòng đất trước sự bàng hoàng của Okayu và Korone, khiến cặp đôi hốt hoảng gọi tên cậu: ``Kuon-senpai!?'' / ``Kuon-senpai à!?''

\begin{dialogue}
  \speak{\textit{Sora}} Đây là cảnh mà Kuon-senpai bắt đầu khởi nghiệp này!
\end{dialogue}

\begin{center}
  \uline{(Có một cái đồng hồ đếm ngược ở góc dưới màn hình. Đồng hồ chỉ về không khi cảnh ở dưới xảy ra.)}
\end{center}

Ngay lập tức, khuôn mặt của Kuon hiện lên trên bề mặt mầm cây. Và từ cái mầm bé xíu đó, một tòa nhà chọc trời bỗng chốc mọc lên như nấm sau mưa, vươn cao vút tận mây xanh.

\begin{figure}[H]
  \centering
  \begin{subfigure}[b]{0.45\textwidth}
    \centering
    \includegraphics[width=8cm]{../../../../Image/Book9/2-12g1.png}
  \end{subfigure}
  \quad
  \begin{subfigure}[b]{0.45\textwidth}
    \centering
    \includegraphics[width=8cm]{../../../../Image/Book9/2-12g2.png}
  \end{subfigure}
  \quad
  \begin{subfigure}[b]{0.45\textwidth}
    \centering
    \includegraphics[width=8cm]{../../../../Image/Book9/2-12h.png}
  \end{subfigure}
  \caption{Sự ra đời kỳ diệu của ``Tòa nhà Kuon'' (thay đầu của Korone bằng đầu của cậu).}
\end{figure}

Sự xuất hiện đột ngột của tòa nhà đưa Korone lên tận đỉnh cao nhất. Okayu đứng dưới đất, mắt sáng rực nhìn lên: ``Công ty Kuon ra đời rồi kìa! Koro-san đứng cao quá chừng luôn !''

Korone vẫy tay rối rít: ``Wai\~{}! Ở trên này nhìn xuống đẹp cực kỳ luôn !''

\begin{dialogue}
  \speak{\textit{Sora}} Công ty Kuon hẳn sẽ làm ăn phát đạt lắm đây.
\end{dialogue}

Nàng Khuyển hào hứng nhảy tưng tưng trên nóc nhà mới, còn nàng Mèo thì tiếp tục múa may ăn mừng phía dưới, mặc kệ sự thật rằng cậu quản lý tội nghiệp giờ đã biến thành móng nhà cho họ [\dots].

\bxh{6}{Cơn tức giận của Shion!}{34}{Ba}    % 3-5

[\dots] Cả Kuon và nàng Hiệp sĩ đều trố mắt kinh ngạc, trong khi nàng Thỏ chỉ biết ngơ ngác nhìn đàn chị đầy khó hiểu. Còn nàng Chó nâu thì nhảy cẫng lên reo hò: ``Đó! Em đã bảo rồi mà! Không sai đi đâu được!''

``Tại sao chị lại bị kẹt trong tường thế này!? Rốt cuộc chuyện này là sao hả, Shion-senpai?'' Noel hốt hoảng hỏi.

``Với cả, mấy cuốn ca-ta-lô kia mắc mớ gì lại nằm ở đó?'' Kuon nhíu mày thắc mắc.

Korone nghiêng đầu nhìn đàn chị, che miệng cười khúc khích: ``Trông dáng vẻ của bà lúc này ngố tàu thật đấy á\~{}''

\begin{dialogue}
  \speak{\textit{Sora}} Chắc cái này mình không cần giải thích thêm nữa nhỉ.
\end{dialogue}

\begin{center}
  \uline{(Có một cái đồng hồ đếm ngược ở góc dưới màn hình.)}
\end{center}

Nàng Phù thủy đảo mắt nhìn cả nhóm, khuôn mặt ửng đỏ vừa xấu hổ vừa cố tỏ ra bình thường thanh lịch: ``À thì\dots\ chị đang hăng say luyện tập phép dịch chuyển tức thời\dots\ nhưng mà hơi trượt tay tính toán sai tọa độ nên là\dots''

Cậu khoanh tay trước ngực, nhíu mày nghiêm nghị chất vấn: ``\dots Thế nên mới kẹt nửa người trong tường như này hả? Thế còn mấy cuốn ca-ta-lô thì sao? Anh không nghĩ chúng nó mọc cánh tự chui vào tường đâu.''

Shion cười gượng gạo, toát mồ hôi hột: ``Ừm\dots\ thì\dots\ em dùng chúng làm vật mẫu để thực hành đó. Ai mà ngờ sức mạnh của em nó lại thành công đến mức\dots\ vượt cả mong đợi như này!''

Kuon vỗ trán cái đốp, thở dài ngao ngán: ``Thành công tới mức kẹt cứng trong tường thế này\dots\ Em làm anh cạn lời luôn rồi đấy, Shion-chan ạ.''

Cậu bước tới, nắm lấy tay áo định kéo nàng phù thủy ra khỏi mớ bê tông lộn xộn, nhưng rồi khựng lại hỏi: ``Mà này, thế quái nào em có thể kẹt được với cái tư thế siêu thực như thế? Cả đầu, cả tay chân đều bị ép sát mép tường thế này á!?''

``Làm sao mà em biết được cơ chứ!'' nàng Phù thủy bật lại gắt gỏng, ``Đến em còn không hiểu nổi tại sao nó lại thành ra bộ dạng thế này nữa!''

Cả nhóm chỉ biết nhìn nhau, nửa buồn cười, nửa bối rối trước sự vụ không thể nào kỳ lạ hơn này.

\begin{center}
  \uline{(Đồng hồ chỉ về không khi cảnh ở dưới xảy ra.)}
\end{center}

``Đúng là đồ ngốc nghếch mà, peko,'' Pekora bĩu môi mỉa mai đàn chị. Đến cậu Tổng cũng phải gật gù đồng tình với cô: ``Peko-chan phán câu này chuẩn đấy. Không cần phải bàn cãi gì thêm.''

Shion ngay lập tức nhăn mặt, giọng pha chút bức xúc tự ái: ``Ý anh là sao hả, Kuon-senpai!? Mà còn em nữa nhé, Pekora-chan, em có tư cách gì mà nói chị hả? Chẳng phải bản thân em cũng đang có mấy củ cà rốt kẹt dính trong tóc đấy thôi?''

\begin{dialogue}
  \speak{\textit{Sora}} Có lẽ nếu Shion-chan không bị mọi người dồn dập trêu chọc thế này, em ấy sẽ không phát hoảng đến như vậy\dots
\end{dialogue}

\img{3-5c}

``Hai cái đó bản chất hoàn toàn khác nhau đó em ạ,'' Kuon thở dài giải thích, ``một cái là do tai nạn phép thuật ngu ngốc, còn một cái là do sở thích cá nhân người ta tự nguyện nhét vào.''

Nàng Thỏ hừ một tiếng tự đắc, cảm thấy bị ``đụng chạm đến lòng tự tôn'', cô liền vểnh tai lên cất lời khoe khoang: ``Kuon-senpai nói chí phải lắm, peko! Không giống như chị Shion-senpai hậu đậu đâu, mấy củ cà rốt của em là tuyệt tác nghệ thuật, có thể gỡ ra gắn vào bất cứ lúc nào tùy theo ý muốn, peko!''

``Thế á? Chém gió vừa thôi. Vậy em thử rút ra đi cho bọn anh xem,'' cậu nhướng mày thách thức, tỏ vẻ không tin.

Pekora hất mặt tự tin, đưa tay lên khéo léo gỡ nhẹ củ cà rốt khỏi lọn tóc đang tết để chứng minh [\dots].

\bxh{5}{Ayame bất ngờ!}{81}{Sáu}    % 6-11

\begin{dialogue}
  \speak{\textit{Sora}} Cái này khỏi phải chờ!
\end{dialogue}

\begin{center}
  \uline{(Chiếc đồng hồ đếm ngược về không ngay lập tức.)}
\end{center}

Cảnh quay bắt đầu với một chiếc ti-vi tĩnh lặng, không có lấy một bóng người, như một bức tranh tĩnh vật đang chờ đợi sự sống động.

Đột nhiên, từ dưới khung hình, Ayame bật lên với khuôn mặt tươi rói như ánh nắng ban mai, reo lên một cách đầy bất ngờ: ``\textbf{Yo-dayo!}''

\gif{ayame1s}{1}{32}

\begin{dialogue}
  \speak{\textit{Sora}} Hẳn là sẽ có một làn sóng tweet ập vào đây.
\end{dialogue}

Cảnh chuyển nhanh sang một góc khác trong phòng, nàng Oni vừa đi vừa nói, giọng đầy hào hứng như một người dẫn chương trình chuyên nghiệp: ``Dạo gần đây trời cũng bắt đầu lạnh rồi nhỉ?''

Đằng sau ống kính, Kuon thầm nghĩ, một nụ cười khẽ hiện trên môi: ``\textit{Thì bây giờ chuẩn bị vào đông rồi mà.}''

Bất thình lình, cô nàng dí sát mặt vào máy quay, đôi mắt ánh lên vẻ tinh nghịch, đến mức có thể nhìn rõ từng chi tiết trên gương mặt cô: ``\textbf{Đúng chứ?}''

Cậu Tổng giật mình, vội vàng giữ chắc máy quay để không bị rung, cảm giác như trái tim mình cũng nhảy theo phản xạ đó.

Cảnh tiếp tục chuyển đến một góc khác, nơi nàng Quỷ sừng đang nằm sấp trên sàn, khẽ lắc lư đầu qua lại với vẻ mặt thoải mái như một chú mèo đang tận hưởng ánh nắng. Giọng cô vẫn tràn đầy năng lượng, lan tỏa khắp căn phòng: ``Mọi người cũng đồng ý với mình chứ?'' [\dots]

\bxh{4}{Tiếng máy cắt cỏ!}{91}{Bảy}   % 7-8

\pov{Tokino Sora xen kẽ với Hikawa Kuon}

\begin{dialogue}
  \speak{\textit{Sora}} Cái này cũng khỏi phải chờ! Cẩn thận nha!
\end{dialogue}

\begin{center}
  \uline{(Chiếc đồng hồ đếm ngược về không ngay lập tức.)}
\end{center}

[\dots] Bên trong văn phòng, tình hình cũng chẳng khá hơn chút nào. Hệ thống sưởi ấm của tòa nhà vẫn chưa được lắp đặt đúng hạn, nên không khí lạnh lẽo từ bên ngoài cứ thế tràn vào, ngưng đọng trong từng góc phòng, làm cho nhiệt độ trong nhà chẳng ấm hơn ngoài trời được bao nhiêu. Những luồng hơi lạnh len lỏi qua từng khe cửa, khe tường, khiến ai cũng phải co ro trong lớp áo khoác dày. [\dots] Thế nhưng, có một người trong căn phòng này dường như chẳng hề bận tâm đến cái lạnh, thậm chí còn đang tạo ra những tiếng động kỳ quặc, như thể cô đang cố gắng tự sưởi ấm cho mình theo một cách nào đó mà không ai hiểu nổi.

``\textbf{Rừm! Rặc! Rặc! Rặc! Rặc! Rặc! Rặc!}'' tiếng của nàng Yêu tinh tuyết - Lamy - vang lên, phá tan cả sự im lặng vốn có.

\gif{7-9i}{1}{82}

Tôi nhìn Lamy - nàng Yêu tinh tuyết - đang vừa rung người, vừa ngước mắt lên trần nhà, miệng không ngừng phát ra những âm thanh khó hiểu. Em ấy ngồi đó với vẻ mặt hờ hững như thể đây là một hoạt động bình thường vào những ngày lạnh giá. ``Em kêu nãy giờ được nửa tiếng rồi đấy. Nãy giờ không thấy chán sao?''

Nàng Yêu tinh tuyết chẳng có vẻ gì là nghe thấy tôi nói, vẫn tiếp tục tạo ra những âm thanh kỳ quặc ấy, đôi mắt cô vẫn hướng về phía trần nhà như đang nhìn một thứ gì đó xa xăm.

Bên cạnh tôi, Ryuuhou cũng đang co ro trên ghế, cô bé nhìn nàng Lamy với vẻ mặt vừa tò mò vừa bối rối: ``Onii-san, Lamy-san đang làm gì thế ạ? Em không hiểu gì cả.''

Tôi chỉ nhún vai với vẻ mặt đầy bất lực, không biết phải giải thích thế nào với em gái.

Nàng Sói đen và nàng Yêu tinh vàng, hai nạn nhân vô tội của tình huống này, đã lặng lẽ trốn ra sau kệ tủ, chỉ ló mỗi cái đầu ra để quan sát từ xa, như thể đang xem một bộ phim kinh dị.

\img{7-9b}

Flare cau mày, giọng cô đầy khó hiểu: ``Lamy-chan đang làm cái quái gì thế?''

Mio liếc đồng đội, giọng thản nhiên như thể đã quen với cảnh tượng này: ``Chắc là đang bắt chước tiếng máy cắt cỏ vì đang chán ấy mà.''

Tôi nhíu mày, quay sang nàng Yêu tinh vàng: ``Lúc buồn chán bộ em ấy toàn làm thế suốt à\dots?''

Nàng Yêu tinh vàng thở dài: ``Có đầy thứ mà em ấy có thể bắt chước được\dots''

Tôi và nàng Sói đen nhìn nhau, cùng đồng thanh chốt hạ với vẻ mặt của những người đã từ bỏ hy vọng: ``\dots mà lại bắt chước cái thứ của nợ đó. Lạ ghê.''

Ryuuhou ngồi bên cạnh, cô bé cũng lắc đầu: ``Em thấy Lamy-san kỳ lạ thật đấy. Mà em thắc mắc là chị ấy có lạnh không nhỉ? Trông chị ấy vẫn thảnh thơi như thường ấy.''

``Em ấy là yêu tinh tuyết xứ tuyết mà,'' tôi nhún vai đáp lại. ``Nhiệt độ từng này đã là gì so với em ấy đâu.'' 

\begin{dialogue}
  \speak{\textit{Sora}} Chính điều này đã làm cho máy cắt cỏ bán đắt hàng như tôm tươi, dẫn tới hủy hoại hết cây xanh trên Trái Đất này\dots
\end{dialogue}

Sau một hồi phát ra những âm thanh vô nghĩa đến mức tôi bắt đầu tự hỏi liệu em ấy có bị ``lỗi hệ thống'' không, Lamy cuối cùng cũng chịu dừng lại. Cô thở phào, lấy ra một chai rượu sake từ đâu đó rồi tự lẩm bẩm với vẻ mặt hài lòng: ``Được rồi, nghỉ một chút nào [\dots].''

{\color{bronze}\bxh{3}{Tiếng kẻng!}{74}{Sáu}}   % 6-5

\pov{Tokino Sora pha lẫn vào ngôi kể thứ ba}

[\dots] Chủ nhân của chiếc kẻng - Aki - xuất hiện từ cánh cửa, bước chân nhẹ nhàng như một linh hồn bay lướt. Cô tiến đến gần Matsuri, người đang cầm chiếc kẻng trong tay với tư thế như một pho tượng cổ, và hỏi với giọng nhẹ nhàng, đầy tò mò: ``Bà có chơi được không, Matsuri-chan?''

Cô nàng Lễ hội không trả lời mà chỉ nở một nụ cười thật tươi, mắt nhắm lại, như thể là một bức tượng phật trong trạng thái hỉ lạc vĩnh cửu; một nụ cười bất diệt đến mức có thể gây rùng mình cho bất kỳ ai nhìn vào.

Cậu Tổng nhìn khuôn mặt hỉ đó, thở dài bất lực, đưa tay xoa trán: ``Em ấy đã như thế khoảng 15 phút rồi.''

Nàng Sói đen, người đã tận mắt chứng kiến toàn bộ quá trình, cúi gằm mặt xuống tỏ vẻ mệt mỏi, đôi vai rũ xuống: ``Tâm của chị ấy trong sáng đến mức mỉm cười 24/7 luôn rồi\dots''

Nàng Dị giới tròn mắt ngạc nhiên, hốt hoảng lui lại một bước: ``Ehh!? Sợ thế!''

Ngồi bên cửa sổ, cô nàng mắt hai màu nhìn Matsuri với vẻ thích thú, thì thầm với anh trai: ``Nhìn chị ấy như một vị phật cười trong chùa ấy. Có khi em nên xin chị ấy một chút `ánh sáng thiền' để áp dụng cho ban nhạc nhỉ.''

\begin{dialogue}
  \speak{\textit{Sora}} Matsuri-chan đã mượn được chiếc kẻng ma thuật từ Aki-chan. Tùy thuộc vào độ thuàn khiết của tâm người đánh mà tiếng của nó sẽ thay đổi!
\end{dialogue}

``Ít nhất thì cái tiếng kẻng chắc phải nghe hay hơn rồi nhỉ?'' cậu Tổng chống hông, nhìn vào chiếc kẻng trong tay Matsuri, giọng đầy hoài nghi.

Nàng Sói đen đứng dậy, đặt tay lên vai Matsuri, rồi hướng ánh mắt về chiếc kẻng trong tay cô nàng, giọng đầy hy vọng: ``Vậy em chơi thử đi!''

Nàng Dị giới hào hứng vỗ tay, như một đứa trẻ sắp được xem ảo thuật: ``Được rồi! Một, hai\dots''

Kuon, Aki và Mio bắt đầu đếm nhịp đồng thanh: ``Hai, ba!''

\begin{center}
  \uline{(Có một cái đồng hồ đếm ngược ở góc dưới màn hình. Đồng hồ chỉ về không khi cảnh ở dưới xảy ra.)}
\end{center}

Nàng Lễ hội nâng chiếc gậy lên, rồi gõ mạnh một cái vào kẻng.

*KENGGGGGGGGGGGGGGGGGGGGGGGGGGGGGGGGGGG!!*

\gif{6-5j}{1}{152}

Matsuri với đôi mắt vẫn nhắm nghiền và nụ cười như pho tượng nâng chiếc gậy lên, rồi gõ mạnh một cái vào kẻng.

*KENGGGGGGGGGGGGGGGGGGGGGGGGGGGGGGGGGGG!! *

\gif{6-5j}{1}{152}

Lập tức, một thứ âm thanh vang vọng cả văn phòng, lan ra khỏi tòa nhà, len lỏi đến thành phố, và rồi cả Trái Đất, như một cơn sóng âm thanh cuốn phăng mọi thứ trên đường đi.

Âm thanh trong trẻo và vang vọng đến mức người viết câu chuyện này ở Hà Nội còn có thể nghe thấy (!?), một minh chứng sống động cho sức mạnh vượt không gian, thời gian của âm thanh.

Cô em gái của cậu Tổng vội bịt tai, nhưng không giấu nổi vẻ kinh ngạc: ``Ôi trời, tiếng này vượt quá tần số của kẻng thường. Rõ ràng nó đang cộng hưởng với một cái gì đó\dots\ nhưng mà là gì nhỉ?''

\ct

Khoảnh khắc đó, nàng Aki và Mio cũng đột nhiên mỉm cười y như Matsuri, một nụ cười vô thức nhưng lại mang theo sự thanh thản lạ thường, như thể vừa chạm vào cõi niết bàn.

Trong khi đó, Rushia và Kuon đứng chôn chân tại chỗ, trố mắt nhìn ba người bọn họ với sự ngỡ ngàng và khó hiểu, như thể vừa chứng kiến một phép thuật ngoài sức tưởng tượng.

\begin{dialogue}
  \speak{\textit{Sora}} Nhóm `Chị em Hút bụi Ion âm' đã trở thành hiện tượng ở thời điểm đó!
\end{dialogue}

Nàng Dị giới mỉm cười rạng rỡ, tuyên bố với một sự phấn khích đến khó tin, giọng như tiếng chuông báo tin vui: ``Và với tiếng kẻng đó, bọn chị đã lập nhóm `Chị em Hút bụi Ion âm' hoàn chỉnh [\dots]!''

{\color{silver}\bxh{2}{Sức mạnh mùa xuân!}{106}{Tám}}   % 8-10

\pov{Tokino Sora xen kẽ với Hikawa Kuon}

[\dots] Cả ba chúng tôi cùng quay đầu về hướng giọng nói phát ra. Một bóng người hiện dần lên từ hư vô, tỏa ra một thứ ánh sáng hồng nhẹ nhàng, mờ ảo. ``Đây là vùng đất của mùa xuân, và ta là nữ thần cai quản nơi đây.''

Giọng nói này nghe quen quá\dots\ Tôi, Ryuuhou và Nene cùng chỉ tay thẳng vào người đó, đồng thanh reo lên: ``À! Miko-chan(-san/-senpai) kìa!''

Cô em gái của tôi cười lớn, chỉ vào đốm sáng hồng quanh người Miko, gọi to để cô ấy nghe rõ: ``Trông chị hay thế! Đang chuẩn bị cho event cosplay gì mới à? Cả ánh sáng phát sáng này còn được chuẩn bị kỹ quá!''

Bóng người tức thì bật dậy, hét lên đầy bực tức: ``\textbf{Này! Đừng có nói toẹt hết ra như vậy chứ, ne!}'' giọng bẹt đặc trưng của Miko lộ rõ hoàn toàn, chẳng còn che giấu được gì nữa.

Tôi nhún vai, đáp lại một cách vô tư: ``Thì giọng của em ai mà không nhận ra. Giống hệt hồi em làm show mini Ai là triệu phú cho Sui-chan ấy, không thể nhầm lẫn được. Lại còn có `ne' ở sau câu nữa, ai lại không biết!''

Nghe tôi nhắc đến sự kiện cũ, mắt em gái tôi lập tức sáng bừng, quay sang nói với Nene hào hứng: ``Đúng rồi Nene-san, chị đó hài cực! Lúc đó chị Miko đóng vai người dẫn chương trình, làm sai cả câu hỏi, khiến cả buổi livestream cười lăn lộn. Em còn cất clip đó trong điện thoại luôn!''

Nene nghe vậy cũng bật cười, tay che miệng: ``Thật á? Nghe hay thế, phải xem lại mới được!''

Ngay sau đó, tiếng đập bàn vang lên thật to, kèm theo tiếng rít khó chịu của micro vọng lại - Miko chắc chắn bị bắt bài nên nổi điên. Tôi vội bịt tai lại, còn Ryuuhou và Nene thì cười đến mức cong người. Cô nàng Vu nữ hậu đậu này đúng là dễ trêu thật.

\begin{dialogue}
  \speak{\textit{Sora}} Trong công cuộc tìm kiếm ``mMùa xuân'' của mình, Nene-chan vô tình ở trong một không gian kỳ lạ bí ẩn\dots
\end{dialogue}

Sau vài phút trấn tĩnh để nguôi cơn giận, Miko-- à không, phải gọi là ``nữ thần mùa xuân'' mới đúng, lại nghiêm giọng tuyên bố, vẫn giữ cái giọng bẹt và âm ``ne'' quen thuộc: ``Vì nhà ngươi có niềm tin mãnh liệt vào mùa xuân, ta sẽ ban cho ngươi sức mạnh của mùa xuân!''

Nene chớp mắt đầy ngạc nhiên, còn tôi thì cau mày định hỏi: ``Khoan, sức mạnh mùa xuân là cái gì--'' tôi chưa kịp nói hết câu.

Cả người cô nàng Hải cẩu đột nhiên tỏa sáng với một thứ ánh sáng hồng rực rỡ, khiến cả hai người còn lại phải nhắm mắt lại một lúc. Khi mở mắt ra, em ấy đang ngạc nhiên nhìn hai bàn tay mình phát sáng, còn Ryuuhou thì há hốc miệng không ngớt, bước lại gần em ấy để ngắm kỹ hơn: ``Trời ơi, đẹp quá Nene-san ơi! Chị trông như công chúa trong truyện cổ tích thật sự, lấp lánh hết!'' 

Rồi cô quay sang tôi, huých vai tôi nhẹ nhàng: ``Onii-san thấy không? Nãy giờ còn nghi ngờ chị ấy không biết gì về mùa xuân, giờ thì tin chưa?''

Tôi chỉ đứng đó với một dấu chấm hỏi to đùng trên đầu: ``\textit{Thế mà lại đúng là em ấy được chọn? Rằng tại sao chứ?}''

``Nữ thần'' Miko để lại lời nhắn cuối cùng trước khi biến mất, giọng vọng lại như từ một nơi xa xôi: ``Hãy dùng sức mạnh đó để cho họ biết thế nào là lễ độ.''

Nene siết tay lại đầy quyết tâm, mắt sáng bừng lên phấn khích, hét to đáp lại: ``Em sẽ không làm chị thất vọng đâu, cảm ơn nhé, Miko-senpai!''

\begin{center}
  \uline{(Có một cái đồng hồ đếm ngược ở góc dưới màn hình.)}
\end{center}

Nhưng ngay lập tức, giọng nói đầy tuyệt vọng của Miko vang lại từ phía xa: ``Đã bảo rồi! Ta là nữ thần mùa xuân cơ mà, ne! Sao lại cứ gọi nhầm tên mãi thế!''

Tôi bật cười lớn, Nene cũng cười theo, còn Ryuuhou thì ôm bụng cười nghiêng ngả, vừa cười vừa nói lắp bắp: ``Miko-san dễ thương quá! Làm thần mà bị bắt bài ngay lập tức, em chưa gặp ai như chị ấy bao giờ!''

Chúng tôi cười một tràng dài, không dứt được trước phản ứng đáng yêu của Miko khi bị bắt tại trận. Rồi, một ánh sáng hồng lại bao trùm lấy cả ba, những dải hoa anh đào rơi xuống, quây thành một vòng tròn xoáy ấm áp quanh chúng tôi, nhẹ nhàng nâng bổng cả ba lên khỏi mặt đất.

\ct

\begin{center}
  \uline{(Đồng hồ chỉ về không khi cảnh ở dưới xảy ra.)}
\end{center}

Khi tôi tỉnh táo lại thì\dots\ cả ba chúng ta đã quay về văn phòng. Ryuuhou bước vào ngay sau tôi, loạng choạng hai bước để giữ thăng bằng, ánh mắt con bé vẫn còn rực rỡ như vừa trải qua một cuộc phiêu lưu khó quên. Cô đảo quanh căn phòng một lượt, rồi thở phào nhẹ nhõm, quay sang tôi nói: ``Thật sự đã về rồi nhỉ? Cái không gian trắng vô tận kia là cái gì mà kỳ lạ thế, em không thể nào lý giải nổi.''

Tôi vừa mới kịp sắp xếp lại suy nghĩ, định trả lời em gái thì có hai giọng nói cùng lúc vang lên từ phía bàn làm việc: ``A, Kuon-senpai về rồi kìa!''

Flare và Kanata cùng ngạc nhiên khi tôi đột nhiên xuất hiện giữa phòng. Ryuuhou đứng yên một chỗ, mắt vẫn dáo dác nhìn ngó mọi ngóc ngách, như thể vẫn chưa tin rằng mình đã thoát khỏi cảnh kỳ lạ kia.

``Anh vừa đi đâu thế?'' nàng Yêu tinh vàng hỏi, đôi mắt to tròn đầy tò mò.

Tôi chỉ đáp gọn: ``Đuổi theo Nene-chan. Tìm được em ấy rồi.''

Thiên sứ định hỏi thêm là Nene đang ở đâu, nhưng chưa kịp mở miệng, một luồng ánh sáng hồng chói lóa bùng phát ngay phía sau tôi, đi kèm giọng nói đầy kiêu hãnh của Nene: ``Hãy nhìn Super-Nene ta đây, giờ đã được ban cho sức mạnh của mùa xuân!''

Nene hiện ra sau màn ánh sáng hồng hào, khẽ lơ lửng trên không trung trước khi đặt chân nhẹ nhàng xuống sàn, hai tay vẫn giữ tư thế tạo dáng đầy ngầu theo cách em ấy nghĩ. Ryuuhou há hốc miệng, vỗ tay một cái thật to: ``Ồ! Chị Nene bay được thật à? Trông tuyệt vời quá chứ!'' hào hứng không kém gì Nene.

Tôi nhíu mày nhìn cô nàng đang khoe khoang kia, trong đầu đầy dấu chấm hỏi. ``Rồi, vậy cái sức mạnh mùa xuân đó có thể làm được gì?''

Chưa kịp để cô nàng Hải cẩu giải thích, tôi đã thấy em ấy đột nhiên giơ cả hai tay lên, bắn ra một tia sáng màu hồng thẳng vào đầu Kanata.

\img{8-10f}

Tôi vội vàng chuẩn bị tinh thần đỡ một trận la hét vì đau đớn\dots\ nhưng trái ngược với dự đoán, Kanata chỉ tròn mắt, tỏ ra cực kỳ ngạc nhiên, tay đưa lên xoa nhẹ trán: ``A! Cơn đau đầu dai dẳng của chị hết rồi này!''

Tôi chớp mắt liên tục, không thể tin vào tai mình. Hả? Cái gì vậy? Ryuuhou cũng không kém phần bất ngờ, nghiêng người sang hỏi tôi: ``E-Em cũng không ngờ là Nene-san có tài chữa bệnh được đấy. Sức mạnh mà Miko-san ban tặng cho chị ấy có thể mạnh tới cỡ nào vậy?''

Nene cũng ngạc nhiên không kém, em ấy nhìn bàn tay mình rồi ngước nhìn Kanata, vẻ mặt không thể tin nổi. Tôi thì toát mồ hôi hột trước tình huống vô lý này: ``V-Vậy cũng được nữa hả?''

Ngay lúc đó, một cơn gió nhẹ bất ngờ thổi qua căn phòng. Nhưng cơn gió này\dots\ lại mang theo một cánh hoa đào, chỉ là nó không phải hoa thật, mà được vẽ thủ công trên giấy, trôi lơ lửng ma mị giữa không trung, xoay tròn như đang nhảy múa. Ryuuhou chỉ vào vật thể lạ đó, gọi cả nhóm cùng xem: ``Mọi người nhìn kìa, cái cánh hoa giấy đó là gì thế? Trôi đi đâu thế?''

\begin{dialogue}
  \speak{\textit{Sora}} Thế là khái niệm ``bốn mùa luôn thay đổi'' kết thúc tại đây.
\end{dialogue}

Flare liền hít một hơi sâu, đôi mắt ánh lên vẻ xúc động, hai tay chắp lại như đang chiêm ngưỡng một kỳ quan: ``Cơn gió xuân đầu tiên kìa!'' giọng em ấy đầy ngưỡng mộ.

Tôi khoanh tay lại, cau mày nhìn cánh hoa đào giấy màu hồng kia lơ lửng: ``Sao anh cảm thấy tình hình này có gì đó sai sai nhỉ?'' tôi thì thầm với Ryuuhou, không giấu nổi sự hoài nghi.

Nene không trả lời, chỉ búng tay một cái, khiến cho cơn gió kỳ lạ kia lượn một vòng quanh phòng rồi biến mất qua khe cửa. Tôi bắt đầu cảm thấy mọi thứ càng lúc càng phi lý. Ryuuhou thì nhìn theo cánh hoa biến mất, quay sang cười với Nene: ``Chơi được đấy Nene-chan, đùa giỡn còn có cả hiệu ứng đặc biệt nữa.''

Thế rồi, tôi chưa kịp suy nghĩ thêm gì, đã nghe Yêu tinh vàng và Thiên sứ cùng nhau lên tiếng, giọng đầy ngưỡng mộ:

``Nene-chan, cảm ơn em. Em là vị cứu tinh đấy.''

``Một nữ thần! Là nữ thần mùa xuân của chúng ta!''

Ryuuhou quay sang nhìn tôi, nhướng mày đầy hoang mang, rồi lại quay sang chọc hai người chị đang đứng đối diện: ``Này hai chị, hai người bị làm sao thế? Sao đột nhiên lại bày ra trò tôn thờ Nene-san vậy?''

Tôi chỉ biết lắc đầu trước tình huống khó hiểu này, buột miệng: ``\dots Có gì đó quá sai luôn.''

Nhưng Kanata và Flare chẳng để tâm đến lời chúng tôi. Họ bắt đầu hô vang: ``Nữ thần! Nữ thần!'' như thể đang tôn thờ một đấng cứu tinh, hai tay vẫn chắp lại trong tư thế thành kính. Ryuuhou không chịu thua, liền bắt chước theo họ, cũng giơ tay hô to một cách đùa cợt: ``Nữ thần! Nữ thần!'' khiến Nene phải quay sang trêu cô.

Ban đầu, Nene chỉ đứng đó, vẻ mặt ngạc nhiên như chưa kịp hiểu chuyện gì đang diễn ra. Nhưng rồi, chỉ sau vài giây, em ấy chớp mắt một cái, và lập tức lộ ra bản chất thích khoe của mình. Nene khoanh tay trước ngực, cười khoái chí, giọng đầy tự mãn: ``Vậy thì quỳ xuống và tôn thờ em đi! Ha!''

Flare và Kanata không chút do dự mà quỳ xuống, hai tay chắp lại như đang cầu nguyện trước một vị thần thực thụ. Ryuuhou đứng cạnh tôi, bật cười lớn trước cảnh tượng vô lý: ``Thiệt tình, hai chị ấy thật sự quỳ xuống kìa! Chị Nene có ma thuật gì mà lừa được hai người vậy?''

Tôi chỉ có thể đứng đó, nhún vai, thở một hơi dài đầy bất lực. ``\textit{Mình không có tuổi để đọ với mấy em ấy rồi\dots}'' tôi thầm nghĩ, đưa tay lên xoa thái dương.

Rồi, tôi khoanh tay lại, lắc đầu lẩm bẩm: ``Rồi chẳng biết `Super-Nene' này làm được trò trống gì không đây\dots''

{\color{gold}\bxh{1}{Mảnh thiên thạch!}{82}{Sáu}}   % S-14

\pov{Tokino Sora pha lẫn vào ngôi kể thứ ba}

[\dots] Nàng Song tính với toàn bộ sức mạnh còn lại, ném mạnh chiếc đầu giả lên trời với tốc độ siêu thanh như một quả bóng chày, một cú ném chứa đựng tất cả sự phẫn nộ và tổn thương.

``{\Large\textbf{HAACHAMA-CHAMA!!}}''

*VÚT!*

\begin{dialogue}
  \speak{\textit{Sora}} Cảnh mà Haato-chan thoát khỏi hang ổ của Subaru-chan và quay trở lại bầy đàn với một tâm trạng bực dọc.
\end{dialogue}

Chiếc đầu bay thẳng về phía chiếc máy tính của Rushiasaurus, kẻ vốn chẳng quan tâm đến trận đấu mà vẫn đang say sưa với trò chơi của mình, hoàn toàn không biết gì về thảm họa đang lao tới.

*RẦM!* *BÙM*

Cú va chạm cực mạnh. Chiếc thùng máy tính phát nổ với một tiếng nổ khủng khiếp, tạo ra một vụ chấn động kinh hoàng hất văng tất cả những con khủng long xung quanh ra xa, khiến chúng ngã nhào trên mặt đất như những con búp bê vải.

``Có vẻ như đến cả những người không liên quan cũng bị vạ lây,'' cậu Tổng tiếp tục bình luận, giọng vẫn giữ vẻ bình thản.

\begin{center}
  \uline{(Có một cái đồng hồ đếm ngược ở góc dưới màn hình.)}
\end{center}

Rushia - vốn dĩ ngồi ngoài cuộc đấu - bàng hoàng nhìn về phía chiếc máy tính của mình. Màn hình hiển thị dòng thông báo lỗi, ứng dụng bị văng ra, toàn màn hình chỉ còn lại một màu xanh lè chết chóc với hàng đống dòng lỗi và một thông báo như đang trêu người cô. Nàng Pháp sư hoảng loạn cực độ, giọng cô vỡ ra.

\begin{center}
  \uline{(Đồng hồ chỉ về không khi cảnh ở dưới xảy ra.)}
\end{center}

*RẦM!*

Cô đập mạnh xuống bàn trong cơn giận dữ, toàn bộ sức lực dồn vào cú đập đó. ``\textbf{Tôi đã làm gì nên vi chứ!? Tại sao lại là tôi!?}''

\img{6-13s}

{\Large \textbf{*RẦM!*} \textbf{*BÙM!*}}

Cú đập thứ hai mạnh đến mức gây ra một vụ nổ kinh hoàng, thổi bay mọi thứ xung quanh như một quả bom nguyên tử vừa phát nổ. Một vùng đất rộng lớn ở Nhật Bản chìm trong ánh sáng trắng chói lòa, đến mức không thể nhìn thấy gì khác ngoài màu trắng tinh khiết. Từ ngoài quỹ đạo Trái đất, có thể thấy rõ một đợt sóng xung kích khổng lồ lan tỏa ra từ điểm nổ, khiến lớp khí quyển trên bầu trời như bị xé toạc, những luồng khí nóng văng ra ngoài vũ trụ kéo theo những mảnh vỡ của mọi công trình trên mặt đất, tạo thành một vệt sáng chói mắt mà ngay cả các kính viễn vọng từ các hành tinh lân cận cũng có thể quan sát được. Ánh sáng dần dịu đi, để lại một vết sẹo mờ trên bề mặt hành tinh.

Khi ánh sáng tan đi, thế giới đã thay đổi hoàn toàn. Mọi thứ đều trở nên khác lạ, nhưng cũng không hẳn là hoàn toàn xa lạ, như thể thời gian đã bị đẩy lùi về một kỷ nguyên khác.

``Và thế là, kỷ Phấn trắng đã đến hồi chấm dứt.'' Fubuki tiếp lời Kuon với một giọng điệu đầy nuối tiếc, như một nhà sử học đang đọc lại những dòng cuối cùng của một thiên anh hùng ca. ``Họ đã tồn tại kể từ sự bùng nổ của kỷ Cambri, nhưng giờ đây, họ đã rơi vào trạng thái ngủ đông.''

\img{6-13t}

``Họ chỉ nổi lên tiếp tục và phát triển từ thời kỳ Đại Tân sinh,'' Kuon tiếp lời cô, kết thúc bình luận của mình với giọng điệu đầy chắc chắn như một nhà tiên tri.

Bề mặt Trái Đất giờ đây chỉ còn lại một bờ cát trắng trải dài đến vô tận, lác đác vài cây dừa đơn độc vươn mình trong gió biển. Những người sống sót xuất hiện từ đống đổ nát và bụi mù, mang theo những vết thương của cuộc chiến.

Sora vẫn đang ngủ say như chết, bình yên đến lạ thường giữa đống hỗn loạn. Rushia ngồi bệt xuống đất cát, tay vẫn còn run rẩy vì những gì vừa xảy ra. Haato bị cắm đầu xuống cát như một cây viết, hai chân giãy giãy trong sự bất lực. Aki, Miko và Choco từ biển bước lên bờ, toàn thân ướt sũng, nước biển chảy dọc theo cơ thể họ. Ở xa xa trên bầu trời, một hình bóng quen thuộc vẫn đang lởn vởn: linh hồn của Subarusaurus, ánh mắt đầy oán hận, nhìn chằm chằm xuống mặt đất như thể đang hăm dọa những kẻ đã tiêu diệt mình, một lời cảnh báo cho tương lai.

\img{6-13u}

\begin{center}
  -Kết thúc.-
\end{center}

\il

Cảnh phim quay lại với bối cảnh tận thế ban đầu.

Gió gào thét như tiếng kêu than của một hành tinh đang hấp hối. Những ngọn lửa bật lên liên tục từ lòng đất, hòa vào bầu trời đỏ lòm như máu đang sôi sục. Giữa đống đổ nát đó, Sora đứng rùng mình giữa những mảnh gạch vụn và màn khói bụi mịt mù.

Nàng Thần tượng cúi nhìn đường nứt to đang lan rộng dưới chân mình: ``Kinh thật! Chỗ này cũng không còn trụ được bao lâu nữa đâu!''

Ngay sau khi nói xong, cô lập tức quay đầu chạy vội, bụi đất tung bay dưới chân như trong một pha hành động bom tấn. Đúng lúc cô vừa kịp la lên: ``\textbf{Hẹn gặp mọi người vào tuần sau!}''

Một thiên thạch khổng lồ rú lên từ vòm trời, lao thẳng xuống đúng vị trí cô vừa đứng. Một cột sáng trắng bệch bao phủ toàn bộ khu vực. Tất cả âm thanh lập tức bị xóa sạch, mọi hình ảnh tan biến, chỉ còn một mảng trắng mênh mông như thể cả thế giới vừa bị xóa vĩnh viễn khỏi bản đồ.

Màn hình dần tối đi\dots
%==========================================TRUYỆN PHỤ=========================================
\omk

\pov{Hikawa Kuon}

\begin{center}
  \uline{Trên xe buýt.}
\end{center}

Trước mặt tôi là một cảnh kết thúc của thước phim, với hình ảnh A-san khoanh tay tạo thành dấu chữ X, gương mặt cười méo xệch, bên cạnh là dòng chú thích đặt ở hai bên: ``Thảm họa của đoàn sản xuất chính là đây.''

\img{s-22c}

Tôi chỉ biết lẩm bẩm tự nhủ: ``\dots Cái gì nữa đang xảy ra thế này?''

Thực ra thì, tôi vừa bị ép phải xem hết một bảng xếp hạng khác, gom chung các đoạn phim cũ của \textit{hologra}, lần này chọn chủ đề là ``ngày tận thế.'' Và cũng như mọi lần, tôi đoán trước được kết quả sẽ ra sao.

Một bảng xếp hạng mà tôi thấy không hề chính xác cho lắm.

Nói thật lòng ra, đối với tôi, ngày nào ở bên đám người kia cũng có thể xem là một phiên bản thu nhỏ của ngày tận thế. Không có buổi làm việc nào trôi qua một cách êm đềm. Cũng chẳng có thứ gì gọi là ``bình thường'' bao giờ. Nếu tôi là người mới vào nghề, chắc đã nộp đơn nghỉ việc từ tuần thứ hai rồi. Nhưng sau hai năm rưỡi sống sót qua đủ trò lằng nhằng của họ, tôi đã trở nên miễn nhiễm với khái niệm ``hỗn loạn.''

Mà nếu nói cho công bằng, cái gọi là ``tận thế'' đúng nghĩa thì tôi xin ``vinh dự'' bổ sung thêm những lần có trò Old Maid. Không chỉ một mà tận hai lần\footnote{Đó là các lần ở {\flaregothic ÉPISODE 99}, quyển số Tám và ở {\flaregothic ÉPISODE 109}, quyển số Chín.}, chỉ với một bộ bài và một lời nói dối nhỏ của ai đó đã khiến chúng tôi phải đền bù gần nửa cái văn phòng, làm tôi phải cấm hẳn trò chơi này mãi mãi. Máy lạnh, máy quay, một nửa hệ thống đèn chiếu, thậm chí cả máy pha cà phê cũng không chẳng thể nào sống sót qua được.

Tôi khoanh tay lại, gật đầu một cách mệt mỏi: ``Nhưng nếu xét `tận thế' theo nhiều nghĩa khác nhau thì\dots\ phải thừa nhận là nó cũng đúng ở một mức độ nào đó. Mà, cái này là fan bình chọn mà, chắc chắn phải có lý do gì đó chứ.''

Tôi thở một hơi dài hơn mức cần thiết. Tôi không biết phải bình luận gì về đoạn phim vừa rồi. Mỗi lần tôi định lên tiếng phản đối điều gì, lại tự nghe câu đó vang vọng trong đầu: ``Đây là \hll\ mà.''

Một câu nói tôi đã lặp đi lặp lại nhiều hơn cả số lần tôi phải trả lời email khẩn từ ban kỹ thuật, và tôi hiểu - không, tôi chấp nhận - là mình không thể nào thoát khỏi chuyện này được nữa.

Tôi còn có lựa chọn nào khác ngoài việc sống chung với đám người đó đâu?

\dots À mà, tôi chưa kịp nói cái phông nền tận thế phía sau Sora lúc dẫn chương trình chỉ là một đoạn CGI miễn phí tôi tải về từ một trang chia sẻ mô hình 3D thôi mà. Không rõ ai đã thêm cái quả thiên thạch vào, nhưng nhìn lại thì\dots\ cũng khá hợp mắt đấy.
\end{document}